\hypertarget{appendix-d49161831368}{}\label{referencestart:appendix-d49161831368}
\section{Bitwise XOR}\label{appendix-d49161831368-bitwise-xor}

\hypertarget{appendix-d49161831368-bitwise-xor}{}

\textbf{Curriculum depth:} Introduced · \textbf{First seen:} C001 · \textbf{Latest development:} C001

\begin{allmightroledefinition}

\subsection{Definition and mechanism}\label{appendix-d49161831368-definition-and-mechanism}

\hypertarget{appendix-d49161831368-definition-and-mechanism}{}

Toggles selected lanes and marks differences; applying the same mask twice restores the input.

\end{allmightroledefinition}

XOR returns 1 when inputs differ. \BookInlineCode{{0}{x}{A}{D}{\ }{\^{}}{\ }{0}{x}{0}{8}{\ }{=}{\ }{0}{x}{A}{5}}{0xAD\ \^{}\ 0x08\ =\ 0xA5} toggles visible; applying \BookInlineCode{{0}{x}{0}{8}}{0x08} again restores \BookInlineCode{{0}{x}{A}{D}}{0xAD}. XOR between old and new words marks changed positions. Reversibility alone does not provide secure encryption.

\subsection{Prerequisites}\label{appendix-d49161831368-prerequisites}

\hypertarget{appendix-d49161831368-prerequisites}{}

\begin{itemize}
\tightlist
\item
  \hyperref[appendix-24475e1596a8]{Boolean algebra} --- The logical domain is false/true.
\end{itemize}

\subsection{Used by}\label{appendix-d49161831368-used-by}

\hypertarget{appendix-d49161831368-used-by}{}

\begin{itemize}
\item
  \href{../01-Bits-Programs-and-Reliable-Systems-P02/main.pdf\#nameddest=appendix-ed338677de0a}{Checksums CRCs and hashes (Bits, Programs, and Reliable Systems, Part 2, Reference notes: Checksums CRCs and hashes)} --- uses this mechanism to define its representation, error or validation contract.
\item
  \hyperref[appendix-1652d27a8e90]{Half adder} --- uses bitwise xor as part of its representation or reasoning.
\end{itemize}

\subsection{Recall and next depth}\label{appendix-d49161831368-recall-and-next-depth}

\hypertarget{appendix-d49161831368-recall-and-next-depth}{}

Explain bitwise xor without the label, work the example above, then state the assumption whose removal would change the result.

\textbf{Pending:} parity, checksums, coding theory, and cryptographic constructions.

\subsection{Source and navigation}\label{appendix-d49161831368-source-and-navigation}

\hypertarget{appendix-d49161831368-source-and-navigation}{}

\hyperref[01-foundations-i-ch001-from-a-signal-to-an-interpretation]{C001 teaching chapter} · \href{../01-Bits-Programs-and-Reliable-Systems-P03/main.pdf\#nameddest=appendix-fd090d982a15}{Part definitions and mechanisms (Bits, Programs, and Reliable Systems, Part 3, Reference notes: Foundations)} · \href{../07-Vision-XR-and-Playable-Worlds-P01/main.pdf\#nameddest=appendix-4ad47bdd269a}{Dependency Map (Vision, XR, and Playable Worlds, Part 1, Reference notes: Dependency Map)}
